% 第十部：本課脈絡中的位置。出處：課程倉 notes/docs/1151/seminar/W02_0921.md、W05_1012_預習.md；
% 1151/seminar/readings/literature-mentions.tsv；各篇文字層 readings/W04、W12、W13/text/。
\chapter{本課脈絡中的位置}

\qaentry{第十部之一}{與前幾週指定文獻的關係}

\begin{qaquestion}
本文與第 2 週 Tremblay、第 4 週 Bickel 與 Fallon、\\
第 5 週 Kelsen 與 Schmitt 各有什麼關係？
\end{qaquestion}

\begin{qafigure}{本課指定文獻與本文的引用關係，依各篇文字層查得。\label{fig:course}}
% 本課指定文獻與本文的引用關係。灰階；本文用強調色。關係直接寫在各框第二行，不用斜放的線上標籤。
\begin{tikzpicture}[qa,wk/.style={nbox,text width=36mm,minimum height=10mm},
    bus/.style={figGray,line width=0.6pt}]
  \node[wk] (tr) at (0,2.55) {W02　Tremblay 2003\\{\scriptsize\color{dimink}未引本文；舉證責任的方向與本文相反（導讀人對照）}};
  \node[wk] (bi) at (0,0.85) {W04　Bickel 1962\\{\scriptsize\color{dimink}本文 1349 引其「對抗多數的難題」}};
  \node[wk] (fa) at (0,-0.85) {W04　Fallon 2018\\{\scriptsize\color{dimink}第 1 章註 25 引本文 1387--1389}};
  \node[wk] (ks) at (0,-2.55) {W05　Kelsen 1929、Schmitt 1931\\{\scriptsize\color{dimink}本文未引；均問誰作終局決定（導讀人對照）}};
  \node[kbox,text width=22mm,minimum height=10mm] (w) at (5.4,0) {W06\\Waldron 2006};
  \node[wk] (gh) at (10.75,0.85) {W12　Ghosh 2010\\{\scriptsize\color{dimink}以本文為主要對話對象之一}};
  \node[wk] (be) at (10.75,-0.85) {W13　Benhabib 2020\\{\scriptsize\color{dimink}討論 Waldron 的政治處境與預先承諾}};
  \coordinate (lb) at (2.95,0);
  \coordinate (rb) at (7.85,0);
  \foreach \n in {tr,bi,fa,ks} \draw[bus] (\n.east) -| (lb);
  \draw[bus] (lb) -- (w.west);
  \foreach \n in {gh,be} \draw[bus] (\n.west) -| (rb);
  \draw[bus] (rb) -- (w.east);
\end{tikzpicture}

\end{qafigure}

\subsection{第 2 週：Tremblay（2003）}

Tremblay 早於本文三年發表。第 2 週的課堂記錄整理其主張：違憲審查不需要特定授權，司法在社會中已經行使有效的強制力並具有政治權威，單是這個事實就構成承認其一般正當性的充分理由；特定正當性命題推定不正當，除非證明有特定授權，一般正當性命題則推定正當，除非能指出某個特定的道德原則使它不被允許。\footcite[525--562]{tremblay2003}

\begin{readernote}
本文第伍部分要求每一個決定程序都能向不同意的 C\textsubscript{n} 說明為何由這些人、以此程序決定\pg{1387}，等於讓違憲審查承擔說明的責任。兩文的舉證責任分配方向相反，可以請聽者比較：Tremblay 從「所有國家機關都受同一原則拘束」起算，Waldron 從「決定程序須有雙方均能接受的正當性理論」起算。
\end{readernote}

\subsection{第 4 週：Bickel 與 Fallon}

本文導論引 Bickel 的「對抗多數的難題」與「脫軌制度」作為既有反對論的代表\pg{1349~及註~7、8}。\footcite{bickel1962} Fallon 的專書第 1 章在註 25 引本文 1387 至 1389 頁，即 C\textsubscript{n} 的兩個問題與立法機關的答覆；\footcite[181 註 25]{fallon2018} Fallon 另於 2008 年以〈The Core of an Uneasy Case for Judicial Review〉為題回應本文。\footcite{fallon2008}

\subsection{第 5 週：Kelsen 與 Schmitt}

本文未引 Kelsen 與 Schmitt。第 5 週的 Schmitt《憲法的守護者》第一章指出，法學上真正的爭點是由誰決定（quis judicabit）；若以一個可裁判一切疑義的憲法法院取代憲法上有權的機關，這個國事法院就會是與國會、帝國總統及帝國政府並列的一個政治機關。\footcite[100--101]{vinx2015} Kelsen 1928 年的報告則從法秩序的位階構造推出憲法法院。\footcite[22--78]{vinx2015}

\begin{readernote}
Schmitt 與 Waldron 都追問終局決定者的身分，理由不同：Schmitt 以司法的界限為據，認為憲法爭議的決定本質上是政治的；Waldron 承認法院有能力審查立法\pg{1363}，問題在於它應否審查、其判斷應否終局，依據是政治平等。兩人主張的決定者也不同，一為帝國總統與政府，一為民選的立法機關。這一則對照屬導讀人所作，原文之間沒有引用。
\end{readernote}

\qaentry{第十部之二}{後續各週如何引用本文}

\begin{qaquestion}
本課後續各週的指定文獻中，哪些引用、回應本文？
\end{qaquestion}

依課程倉的互見計數，後續指定文獻提及 Waldron 之姓最多者為第 12 週的 Ghosh（2010）與第 13 週的 Benhabib（2020）。

\begin{itemize}
\item 第 12 週 Ghosh（課綱作 Lilian Castro，刊本作者為 Eric Ghosh）以本文為主要的對話對象之一：刊本 330 至 331 頁摘述本文的程序面與結果面理由，引本文 1372、1373 頁；其後回應本文的程序面論證與結果面論證。\footcite{ghosh2010}
\item 第 13 週 Benhabib 討論 Waldron 的「政治的處境」（circumstances of politics）、「立法的尊嚴」與「參與式多數決」，並轉述 Waldron 對預先承諾之說的反對。\footcite[509--513]{benhabib2020}
\end{itemize}
